\chapter{Wzorce uzycia i gotowe przepisy}

\section{Po co w ogole wzorce}
Gdy poznasz juz skladnie jezyka, kolejnym krokiem nie jest zapamietanie kolejnych tokenow,
lecz nauczenie sie powtarzalnych sposobow rozwiazywania realnych problemow. Wzorzec to
po prostu sprawdzony uklad kodu, ktory warto miec z tylu glowy.

\section{Wzorzec: strona glowna z sekcjami}
Bardzo czesty przypadek to jedna strona glowna z kilkoma sekcjami. W MoonChunk warto
podzielic ja na przygotowanie danych i kilka blokow \mc{content}.

\begin{lstlisting}[style=moonchunk,caption={Strona glowna skladana z sekcji}]
chunk "Home" {
  output: "./dist";
  title: "Strona glowna";

  let heroTitle: string = "MoonChunk";
  let heroText: string = "Jezyk do generowania statycznych stron.";

  page "" {
    content {
      <section class="hero">
        <h1>{heroTitle}</h1>
        <p>{heroText}</p>
      </section>
    };

    content {
      <section class="cta">
        <a href="/docs">Przejdz do dokumentacji</a>
      </section>
    };
  };
};
\end{lstlisting}

Ten wzorzec jest prosty, czytelny i dobrze skaluje sie do landing page'y.

\section{Wzorzec: lista kart}
Jeśli chcesz zbudowac siatke podobnych elementow, np. kart, petla jest naturalnym wyborem.

\begin{lstlisting}[style=moonchunk,caption={Lista kart budowana z tablicy obiektow}]
const cards = [
  { title: "Parser", description: "Rozpoznawanie skladni." },
  { title: "Runtime", description: "Wykonywanie programu." },
  { title: "Output", description: "Generowanie HTML." }
];

page "" {
  for (let i: int = 0; i < 3; i++) {
    content {
      <article class="card">
        <h3>{cards[i].title}</h3>
        <p>{cards[i].description}</p>
      </article>
    };
  };
};
\end{lstlisting}

Warto zauwazyc, ze najwygodniej przechowywac dane kart jako obiekty, a nie jako kilka
osobnych tablic.

\section{Wzorzec: stan pusty}
To bardzo wazny scenariusz w praktyce. Jesli lista jest pusta, strona nie powinna byc
po prostu biala i milczaca. Lepiej pokazac użytkownikowi jasny komunikat.

\begin{lstlisting}[style=moonchunk,caption={Pusty stan listy}]
let hasItems: bool = false;

page "" {
  if (hasItems) {
    content {
      <p>Mamy elementy do pokazania.</p>
    };
  } else {
    content {
      <p>Na razie brak elementow.</p>
    };
  };
};
\end{lstlisting}

\section{Wzorzec: filtr publikacji}
W projektach tresciowych bardzo czesto przechowujesz wpisy, ktore nie powinny jeszcze
trafiac do publicznego HTML-a. Najprostszy wzorzec wygląda tak:

\begin{lstlisting}[style=moonchunk,caption={Pokazuj tylko elementy opublikowane}]
for (let i: int = 0; i < 3; i++) {
  if (posts[i].published) {
    content {
      <h2>{posts[i].title}</h2>
    };
  };
};
\end{lstlisting}

To nie tylko czytelne, ale i bezpieczne. Szkice nie wyplywaja przypadkiem do finalnej
strony.

\section{Wzorzec: funkcja do etykiety}
Jesli wiele miejsc potrzebuje podobnej etykiety, nie powtarzaj warunku w kółko.

\begin{lstlisting}[style=moonchunk,caption={Funkcja etykietujaca stan}]
function statusLabel(published: bool): string {
  return published ? "Opublikowany" : "Szkic";
}
\end{lstlisting}

Taki wzorzec bardzo dobrze dziala dla statusow, podpisow i krotkich komunikatow.

\section{Wzorzec: wspolne metadata plus lokalne nadpisanie}
To jeden z najbardziej produktywnych wzorcow w MoonChunk.

\begin{lstlisting}[style=moonchunk,caption={Wspolne metadata i lokalne doprecyzowanie}]
export chunk "Metadata" {
  lang: "pl";
  charset: "utf-8";
  metaAuthor: "MoonChunk Team";
};

chunk "About" {
  @include Metadata;
  output: "./dist";
  title: "O projekcie";
  metaDescription: "Opis strony o projekcie.";
};
\end{lstlisting}

Dzieki temu projekt ma centrum konfiguracji, ale pojedyncze strony nadal zachowuja
swoja tozsamosc.

\section{Wzorzec: oblicz najpierw, pokaz potem}
To wzorzec, ktory warto stosowac niemal wszedzie.

\begin{lstlisting}[style=moonchunk,caption={Oddzielenie logiki od prezentacji}]
let canShowPromo: bool = views > 100 and isPublished;
let promoText: string = canShowPromo ? "Popularny wpis" : "Standardowy wpis";

content {
  <p>{promoText}</p>
};
\end{lstlisting}

Dzieki temu kod HTML nie jest przeciazony zlozonymi wyrazeniami.

\section{Wzorzec: rekurencyjny spacer po etapach}
Rekurencja nie sluzy tylko do klasycznej silni. Mozna jej uzyc tez do przechodzenia po
kolejnych krokach, np. numerowanych etapach procesu.

\begin{lstlisting}[style=moonchunk,caption={Rekurencja jako przechodzenie po krokach}]
function walk(index: int, max: int): int {
  print("step:", index);
  return index >= max ? index : walk(index + 1, max);
}
\end{lstlisting}

To dobry wzorzec wszedzie tam, gdzie problem naturalnie rozklada sie na podobne kroki.
